\subsection{Environment-Dependent Hubble Measurements}
  \label{subsec:h0_environment}

  A primary prediction of the framework is a correlation between the locally inferred
  Hubble constant and the surrounding large-scale environment.

  Observers and standard candles located in or near deep cosmic voids are expected
  to yield systematically higher inferred values of $H_0$ than those located in denser
  regions.
  Conversely, measurements anchored in overdense environments should converge
  more closely toward the global expansion rate $H_{\mathrm{global}}$.

  This prediction can be tested by cross-correlating existing and forthcoming
  $H_0$ measurements with void catalogs and density reconstructions derived from
  large-scale galaxy surveys, as already explored in the context of local
  underdensity scenarios~\cite{Keenan2013KBC,Shanks2019Void,Kennworthy2019VoidCritique}.
  A statistically significant null result would strongly disfavor the effective
  saturation mechanism.

  In particular, the framework predicts a differential signal:
  measurements performed in environments of comparable redshift but contrasting
  large-scale density should exhibit systematically different inferred values of $H_0$.
